\chapter{The Magic Square}

\section{What to do}
You get a \(3\times 3\) empty grid and nine little cards labelled
\(1,2,3,4,5,6,7,8,9\). Place each number in its own cell so that every
\emph{row}, every \emph{column}, and the two \emph{diagonals} add up to the
\emph{same} total. That matching total is the ``magic sum.''

\begin{center}
\begin{tabular}{|c|c|c|}
\hline
\phantom{0} & \phantom{0} & \phantom{0} \\
\hline
\phantom{0} & \phantom{0} & \phantom{0} \\
\hline
\phantom{0} & \phantom{0} & \phantom{0} \\
\hline
\end{tabular}
\end{center}

\noindent
Try it with pencil and paper before you read on---it is a good puzzle for a
rainy afternoon, and you can check your answer by adding the three lines through
the middle in more than one way.

\section{Two puzzles for you}
\begin{enumerate}
  \item Which number \emph{must} sit in the very centre cell?
  \item What is the magic sum---the number that every row, every column, and
  each diagonal should add to?
\end{enumerate}

The next sections think about these questions slowly, using symmetry and
fairness, not letter-maths.

\section{Thinking step by step: the middle number}
Picture the numbers \(1\) through \(9\) written in order. The smallest and the
largest can be paired as friends that balance: \(1\) with \(9\), \(2\) with
\(8\), \(3\) with \(7\), \(4\) with \(6\). Each pair has the same feel: two
numbers that sit the same distance from the middle of the list. The only number
left without a partner is the one \emph{in the middle} of the list---the fifth
number when you count \(1,2,3,4,5,6,7,8,9\)---and that is~\(5\).

Now look at the empty square. The centre cell is special because four different
lines cross there: the middle row, the middle column, and both diagonals. Those
lines should all behave the same way if the square is to stay balanced when you
picture turning it or flipping it. The number that belongs in such a ``balancing
seat'' is not a tiny number stuck on an edge and not a huge number pushed into
a corner; it is the number that already sits in the middle of the pack
\(1,\ldots,9\). So, by symmetry, the centre should hold~\(5\).

\section{Thinking step by step: the magic sum}
Every number from \(1\) to \(9\) is used exactly once. If you add them all
together, you get \(45\). (You can check by pairing \(1\) with \(9\), \(2\) with
\(8\), and so on---each pair makes \(10\)---and then the middle \(5\) is left
over.)

Now think about the three \emph{rows}. They fill the whole square without overlap,
so the three row-totals, taken together, use every number once and therefore add
up to~\(45\). The three rows are supposed to be \emph{equal}, so each row must
get the same share of that \(45\). Picture three piggy-banks that should hold
matching amounts: only \(15\), \(15\), and \(15\) work.

Here is another symmetry picture once you know the centre holds~\(5\). Any line
through the centre (a row, a column, or a diagonal) has two other cells. Those
two numbers should match the pattern of the pairs above: they ought to add to
\(10\) so that, together with the middle \(5\), the line adds to \(15\). So the
magic sum is \(15\).

The same \(15\) works for the columns and diagonals when the square is built
fairly.

\section{The final answer}
Here is the answer to the two puzzles, all in one place:

\begin{center}
\renewcommand{\arraystretch}{1.3}
\begin{tabular}{|c|}
\hline
\textbf{The middle cell must be \(5\).}\\[0.35em]
\textbf{Every row, column, and diagonal adds to \(15\).}\\
\hline
\end{tabular}
\end{center}

\noindent
One completed square that works is below (you can rotate or flip the pattern and
still get a valid magic square):

\begin{center}
\begin{tabular}{|c|c|c|}
\hline
8 & 1 & 6 \\
\hline
3 & 5 & 7 \\
\hline
4 & 9 & 2 \\
\hline
\end{tabular}
\end{center}

\noindent
Check: each line through the grid adds to \(15\), and \(5\) sits in the centre.
Have fun building your own version!

\section{Number Scrabble (Pick 15)}
\emph{Number Scrabble} is also called \emph{Pick~15} or ``3~to~15.'' The rules
are simple:

\begin{itemize}
  \item Two players take turns.
  \item On each turn, a player chooses one of the numbers \(1,2,\ldots,9\) that
  has \emph{not} been chosen before.
  \item As soon as one player owns \emph{three} numbers whose sum is exactly
  \(15\), that player wins.
  \item If all nine numbers have been taken and nobody has a winning triple,
  the game is a draw.
\end{itemize}

\noindent
A readable overview appears on Wikipedia under ``Number Scrabble''; see
\url{https://en.wikipedia.org/wiki/Number_Scrabble}.

\subsection*{Why this is the same game as tic-tac-toe}
Write the numbers \(1\)--\(9\) into a \(3\times 3\) magic square---for instance
the one we already met:

\begin{center}
\begin{tabular}{|c|c|c|}
\hline
8 & 1 & 6 \\
\hline
3 & 5 & 7 \\
\hline
4 & 9 & 2 \\
\hline
\end{tabular}
\end{center}

\noindent
Every row, column, and diagonal adds to \(15\). So a ``line'' of three cells in
this picture is exactly a set of three numbers that sum to \(15\), and nothing
else does. Choosing a number in Number Scrabble is the same as putting your mark
in the matching cell of this square; getting three numbers that sum to \(15\)
is the same as getting three cells in a row, column, or diagonal. In other
words, Number Scrabble and tic-tac-toe are the same game in disguise---only the
drawing looks different.

\subsection*{Optimum play}
Because the games match, every good idea for tic-tac-toe becomes a good idea here.
The first player should usually take~\(\mathbf{5}\) (the centre of the square),
since that number appears in the most winning triples. The second player should
answer in a \emph{corner} number (\(2\), \(4\), \(6\), or \(8\)) rather than an
\emph{edge} number (\(1\), \(3\), \(7\), or \(9\)), for the same geometric
reason as in tic-tac-toe. If both sides keep playing as well as they can, neither
player can force a win: the game ends in a \emph{draw}. That is the optimum
outcome when nobody slips.
